% ============================================================
% appA.tex —— 附录 A 李论基础
% 底本：LNM 2386，PDF p461–p476（印刷 pp.421–436）
% 体例：\chapter* 不编号；节/小节手写编号 A.1 / A.1.1；公式 \tag{A.x}
% ============================================================

% [原文 p.421]
\chapter*{附录 A\quad 李论基础}
\addcontentsline{toc}{chapter}{附录 A\quad 李论基础}
\setcounter{section}{0}
\renewcommand{\thesection}{A.\arabic{section}}
\renewcommand{\thesubsection}{\thesection.\arabic{subsection}}

\ChapterAbstract{在本附录中，为方便读者，我们汇集李论中对理解本书正文所涉及的较深入的李论内容有帮助的基本概念与结果。我们特别聚焦于复约化李群的分类、根系、实李群、Levi 分解与抛物子代数。例如，用 Dynkin 图给出的分类以及各种情形下根的描述，在第 12 章引入李代数上的正性结构与正表示这两个概念时是有用的。进一步阅读的标准综合参考文献包括 [195] 或 [234]。}

\section{复李群及其分类}

\subsection{复约化李代数}

\begin{definition}[A.1] 设 $\g$ 是复有限维李代数。若对 $\g$ 中的每个理想 $\mathfrak{a}$，都存在 $\g$ 中的理想 $\mathfrak{b}$ 使得 $\g=\mathfrak{a}\oplus\mathfrak{b}$，则称 $\g$ 是约化的（reductive）。
\end{definition}

$\g$ 的根基（radical）记作 $\operatorname{rad}\g$，定义为 $\g$ 中极大的可解理想。

复约化李代数 $\g$ 的中心同构于其根基 $\operatorname{rad}\g$。此外有分解
\[
\g=\operatorname{rad}\g\oplus[\g,\g],
\]
其中 $[\g,\g]$ 是 $\g$ 的半单部分。

复半单李代数 $\g$ 的伴随表示是由李括号
\[
\operatorname{ad}_x(y)=[x,y]
\]
定义的李代数表示 $\mathrm{ad}:\g\to\End(\g)$。

存在极大阿贝尔子代数 $\mathfrak{t}\subset\g$，使得伴随作用在 $\g$ 上是对角的。这指的是伴随表示 $\mathrm{ad}|_{\mathfrak{t}}$ 把李

% [原文 p.422]
代数 $\g$ 分解为不可约表示的直和，这些表示称为权空间。对应于权零的子代数 $\mathfrak{t}=\g_0$ 称为 $\g$ 的 Cartan 子代数；任意两个 Cartan 子代数都由 $\g$ 的某个自同构共轭。因此所有 Cartan 子代数的维数相同；我们定义复半单李代数 $\g$ 的秩为该维数的公共值：
\[
\operatorname{rk}_{\C}(\g)=\dim(\mathfrak{t}).
\]

现在设 $\mathfrak{t}\subset\g$ 是 $\g$ 的一个极大阿贝尔子代数（模去 $\g$ 的内自同构）。注意 $\mathfrak{t}\supset\operatorname{rad}\g$。伴随表示
\[
\mathrm{ad}:\g\to\End(\g)
\]
限制到 $\mathfrak{t}$ 上可对角化，于是得到分解
\[
\g=\mathfrak{t}\oplus\bigoplus_{\alpha\in R}\g_\alpha,
\]
其中 $R\subset\mathfrak{t}^{*}$ 是 $\g$ 的根集。对每个 $\alpha$，空间 $\g_\alpha$ 是根 $\alpha\in R$ 的特征子空间，即 $X\in\g_\alpha$ 当且仅当
\[
[h,X]=\alpha(h)X
\]
对每个 $h\in\mathfrak{t}$ 成立。

Killing 形式 $K:\g\times\g\to\C$ 定义为 $\g$ 上的复双线性形式
\[
K(X,Y)=\Tr(\mathrm{ad}X,\mathrm{ad}Y).
\]
当限制到 $\g$ 的半单部分 $[\g,\g]$ 上时，Killing 形式非退化，因此可以引入：

\begin{definition}[A.2] 设 $\g$ 是复约化李群（原文如此；按上下文应为李代数），$\mathfrak{t}\subset\g$ 是极大阿贝尔子代数。对每个根 $\alpha\in\mathfrak{t}^{*}$，余根 $h_\alpha\in\mathfrak{t}$ 是 $[\g,\g]\cap\mathfrak{t}$ 中满足
\[
\alpha=2\,\frac{K(-,h_\alpha)}{K(h_\alpha,h_\alpha)}
\]
的唯一元素。我们用 $C\subset\mathfrak{t}$ 表示 $\g$ 的所有余根之集。
\end{definition}

% [Tutor Note: 定义 A.2 的原文写作「设 $\g$ 是复约化李群」，但下文 $\mathfrak{t}\subset\g$、$\alpha\in\mathfrak{t}^{*}$ 以及余根 $h_\alpha\in\mathfrak{t}$ 都表明 $\g$ 应为李代数（与定义 A.1 的对象一致）。这是原文的一处笔误；译文照录并在括号中注明。]

\begin{remark}[A.3] 余根 $C$ 张成 $\mathfrak{t}$（作为复向量空间）当且仅当 $\g=[\g,\g]$，即 $\g$ 是半单复李代数。
\end{remark}

权空间分解的一个推论是：对线性函数 $l:\mathfrak{t}^{*}\to\R$ 的选取，只要 $l(\Delta)\ne0$，根集 $\Delta$ 就分裂为正根与负根，即
\[
\Delta=\Delta_+\cup\Delta_-,
\]
其中 $\Delta_+:=\{\alpha\in\Delta\mid l(\alpha)>0\}$，而 $\Delta_-$ 类似地定义。我们称正根 $\alpha\in\Delta_+$ 是单的（simple），若它不能写成

% [原文 p.423]
$\Delta_+$ 中元素的正系数线性组合。单根集记作 $\Pi=\{\alpha_1,\dots,\alpha_n\}$，其中 $n$ 是 $\g$ 的秩，它构成李代数 $\mathfrak{t}^{*}$ 的一组基。

若 $\mathrm{ad}_e$ 作为 $\g$ 上的线性映射是幂零的，则称元素 $e\in\g$ 是幂零的。由 Morozov–Jacobson 定理，任何幂零元素都嵌入 $\g$ 的一个三维单子代数，即一个同构于李代数 $\sl_2$ 的子代数（证明见 [61, 第 VIII 章, 第 11 节, 命题 2]）。事实上存在另一个幂零元素 $\tilde e\in\g$，使得
\begin{equation}
\langle e,\tilde e,h\rangle\cong\sl_2(\C), \tag{A.1}
\end{equation}
并满足关系 $[h,e]=e$、$[h,\tilde e]=-\tilde e$、$[e,\tilde e]=h$。特别地，若
\[
e=\sum_{\alpha\in\Delta_+}c_\alpha e_\alpha,
\]
其中 $c_\alpha$ 为正整数、$e_\alpha\in\g_\alpha$、$\tilde e_\alpha\in\g_{-\alpha}$（$\alpha\in\Delta_+$），则称该幂零元素为正则幂零或主幂零。

\begin{definition}[A.4] 复半单李代数 $\g$ 的一个子代数称为主 $\sl_2$-子代数，若它同构于 (A.1) 中的 $\sl_2$，且含有 $\g$ 的正则幂零元素 $e$。这类子代数只有一个共轭类，因此我们有时直接说「$\g$ 的主 $\sl_2$-子代数」。
\end{definition}

下面我们定义：

\begin{definition}[A.5] 设 $\g$ 是复半单李代数，$\mathfrak{t}$ 是 $\g$ 的 Cartan 子代数，$\Delta_+$ 是正根集。相对于 Cartan 子代数与正根集的选取，标准 Borel 子代数定义为子代数
\[
\mathfrak{b}=\mathfrak{t}\oplus\bigoplus_{\alpha\in\Delta_+}\g_\alpha.
\]
标准 Borel 子代数也称为极小抛物子代数；任何包含 $\mathfrak{b}$ 的子代数称为标准抛物子代数。
\end{definition}

% [Tutor Note: (A.1) 中的 $\langle e,\tilde e,h\rangle$ 表示由这三个元素张成的子代数（而非内积）；它与 $\sl_2(\C)$ 的同构由关系 $[h,e]=e$、$[h,\tilde e]=-\tilde e$、$[e,\tilde e]=h$ 给出，即 $e$ 为升算子、$\tilde e$ 为降算子、$h$ 为 Cartan 元素。核对方式：在 $\sl_2(\C)$ 的标准基 $e=\begin{pmatrix}0&1\\0&0\end{pmatrix}$、$\tilde e=\begin{pmatrix}0&0\\1&0\end{pmatrix}$、$h=\begin{pmatrix}1&0\\0&-1\end{pmatrix}$ 上直接计算三个括号即可。]

\subsection{特征与余特征}

与我们对复李代数给出的定义类似，我们有：

\begin{definition}[A.6] 复李群 $G$ 称为约化的，若它连通且其根基是环面，即根基是同构于乘法群 $\C^{\times}=\C\setminus\{0\}$ 的若干复制的直积的复李群。
\end{definition}

对于复约化李群 $G$ 及其中的极大环面 $T\subset G$，$T$ 对应的李代数 $\mathfrak{t}$ 在 $\g$ 中极大阿贝尔。因此存在

% [原文 p.424]
相应的根集 $R\subset\mathfrak{t}^{*}$ 与相对于 $T$ 的余根集 $C\subset\mathfrak{t}$。我们有：

\begin{definition}[A.7] 对上述 $G$ 与 $T$，$T\subset G$ 的特征格定义为复李群之间态射构成的自由阿贝尔群
\[
\Lambda^{\perp}=\Hom(T,\C^{\times}).
\]
\end{definition}

指数映射 $\exp:\C\to\C^{\times}$，$z\mapsto\ee^z$，给出 $\C^{\times}$ 的李代数与 $\C$ 之间的标准等同。于是有包含
\[
\Lambda^{\perp}:=\Hom(T,\C^{\times})\subset\Hom(\mathfrak{t},\C)=\mathfrak{t}^{*}.
\]
于是按定义有 $R\subset\Lambda^{\perp}$，从而有：

\begin{definition}[A.8] 对上述 $G$ 与 $T$，$T\subset G$ 的余特征格定义为自由阿贝尔群
\[
\Lambda=\Hom(\C^{\times},T)\subset\Hom(\C,\mathfrak{t})=\mathfrak{t}.
\]
\end{definition}

注意 $\g$ 的所有余根都落在余特征格中，即 $C\subset\Lambda$。

回忆向量空间 $\mathfrak{t}$ 与其对偶 $\mathfrak{t}^{*}$ 之间的自然配对。格 $\Lambda$ 与 $\Lambda^{\perp}$ 之间有典范配对
\begin{equation}
\langle\cdot,\cdot\rangle:\Lambda\times\Lambda^{\perp}\to\Z, \tag{A.2}
\end{equation}
它由复合
\[
\Hom(\C^{\times},T)\times\Hom(T,\C^{\times})\to\Hom(\C^{\times},\C^{\times})=\Z
\]
定义，并与向量空间与其对偶之间的配对一致。配对 (A.2) 是完美的，因此 $\Lambda^{\perp}$ 是 $\Lambda$ 的对偶格。

还要注意，把 $\Lambda$ 与 $\Q$ 张量后，配对 (A.2) 可以延拓到 $\Q$ 上。我们称 $\Lambda\otimes_{\Z}\Q$ 中的元素为（有理）权。在微分之下，$\theta$ 可以看作极大环面 $T$ 的李代数 $\mathfrak{t}$ 中的点。类似地，我们称 $\Lambda^{\perp}\otimes_{\Z}\Q$ 中的元素为（有理）余权。

\begin{remark}[A.9] 设 $\{t_i\}$ 是格 $\Lambda$ 的生成元。滥用记号，我们可以把 $\{t_i\}$ 看作李代数 $\mathfrak{t}$ 的一组基。于是有理权 $\theta$（看作 $\mathfrak{t}$ 中相应的元素）可以写成 $\theta=\sum\frac{a_i}{d_i}t_i$。
\end{remark}

\subsection{根数据及其 Langlands 对偶}

在引入特征格与余特征格之后，我们现在有：

% [原文 p.425]
\begin{definition}[A.10] 设 $G$ 是连通复李群，并选定极大环面 $T\subset G$。$G$ 的根数据是如下数据：
\[
\mathcal{R}(G,T):=\{C\subset\Lambda,\ R\subset\Lambda^{\perp}\},
\]
它还隐含地诱导出配对 (A.2) 以及根与余根之间的双射 $R\to C$，$\alpha\mapsto h_\alpha$。
\end{definition}

根数据的概念由 Demazure [104] 引入，作为分类复约化李群的工具。确实，根数据把约化李群 $G$ 确定到同构，这由 Steinberg [334] 证明。

对于复约化李群 $G$ 及其极大环面 $T\subset G$ 的选取，我们可以交换根与余根以及特征格与余特征格的角色，从而得到新的根数据
\[
\mathcal{R}(G,T)^{\vee}:=\{R\subset\Lambda^{\perp},\ C\subset\Lambda\},
\]
称为对偶根数据。这引出：

\begin{definition}[A.11] 对上述 $G$ 与 $T$，$G$ 的 Langlands 对偶群是由对偶根数据 $\mathcal{R}(G,T)^{\vee}$ 确定的（在同构意义下唯一的）复约化群。我们用 $G^{\vee}$ 表示 $G$ 的 Langlands 对偶群。
\end{definition}

\begin{example}[A.12] $\PSL_2(\C)$ 的 Langlands 对偶是 $\PSL_2(\C)^{\vee}=\SL_2(\C)$。更一般地，无中心约化群的 Langlands 对偶是单连通的。
\end{example}

% [Tutor Note: 例 A.12 原文印作「The Langland dual of $\PSL_2(\C)$」，专名漏了末尾的 s（应为 Langlands）；译文按正确的专名「Langlands 对偶」处理，特此说明。另核对：$\PSL_2(\C)$ 无中心，故其对偶单连通，确为 $\SL_2(\C)$。]

\subsection{根系}

有限维向量空间的根系概念对数学的许多领域都很重要，但在单李代数的研究中尤为本质，并导致这些李代数的分类。我们在此回顾其一些基本特征，关于非常透彻的讨论可参见 [195, 第 10 章, 第 3 节]。

\begin{definition}[A.13] 设 $V$ 是有限维实向量空间，$\alpha$ 是 $V$ 中非零元素。沿 $\alpha$ 的反射是 $V$ 上满足以下条件的线性变换 $s$：
\begin{enumerate}
\item $s(\alpha)=-\alpha$；
\item $s$ 的不动点构成 $V$ 中的一个超平面。
\end{enumerate}
\end{definition}

注意任何反射的阶为 2，且由不动点确定。反射的一个重要性质是：给定非零元素 $\alpha\in V$，至多存在一个沿 $\alpha$ 的反射使某个有限生成子集 $R\subset V$ 不变。我们现在引入：

\begin{definition}[A.14] 设 $V$ 是有限维实向量空间，$R\subset V$ 是有限个非零向量组成的集合。若以下条件满足，则称 $R$ 是 $V$ 中的根系，其元素称为根：

% [原文 p.426]
\begin{enumerate}
\item 集合 $R$ 生成向量空间 $V$。
\item 对每个 $\alpha\in R$ 存在沿 $\alpha$ 的反射 $s_\alpha$ 使 $R$ 不变；此反射必须是唯一的。
\item 对所有 $\alpha,\beta\in R$，由方程 $s_\alpha(\beta)=\beta-a_{\alpha,\beta}\alpha$ 确定的数 $a_{\alpha,\beta}$ 是整数。
\end{enumerate}
\end{definition}

若对 $R$ 中任一对元素 $\alpha,\beta$，当 $\beta=m\alpha$（$m$ 为实数）时必有 $m=\pm1$，则称该根系是约化的。

现在设 $\g$ 是复单李代数，$\mathfrak{t}$ 是 Cartan 子代数，即 $\g$ 的自正规化幂零子代数。$\g$ 相对于 $\mathfrak{t}$ 的根是非零元素 $\alpha\in\mathfrak{t}^{*}$，使得存在非零 $X\in\g$ 满足
\[
[H,X]=\alpha(H)X
\]
对每个 $H\in\mathfrak{t}$ 成立。$\g$ 相对于 Cartan 子代数 $\mathfrak{t}$ 的根集 $\Delta=\Delta(\g,\mathfrak{t})\subset\mathfrak{t}^{*}$ 是有限的，并且按定义 A.14 是约化根系。

对 $\mathfrak{t}^{*}$ 中的根系 $R$，用 $\Aut(R)$ 表示 $\mathfrak{t}^{*}$ 的保持 $R$ 不变的线性变换构成的有限群。$\Aut(R)$ 中由反射 $\{s_\alpha\mid\alpha\in R\}$ 生成的子群 $W(R)$ 称为 $R$ 的 Weyl 群。我们进一步定义子集 $B\subset R$ 为根系 $R$ 的一组基，若：
\begin{enumerate}
\item 集合 $B$ 是 $\mathfrak{t}^{*}$ 的一组基；
\item 每个元素 $\beta\in R$ 都可写成 $\beta=\sum_{\alpha\in B}n_\alpha\alpha$，其中系数 $n_\alpha$ 为同号的整数。
\end{enumerate}
复单李代数的每个根系都有基，且任意两组基都由唯一的 Weyl 群元素共轭。

存在唯一的正定内积 $\langle,\rangle$ 在 $\mathfrak{t}^{*}$ 上、在 $W(R)$ 下不变，满足
\[
\langle\lambda,\mu\rangle=\sum_{\alpha\in R}\langle\lambda,\alpha\rangle\langle\alpha,\mu\rangle,
\qquad\lambda,\mu\in\mathfrak{t}^{*}.
\]
对定义 A.14 中的 $a_{\alpha,\beta}$，于是有
\[
a_{\alpha,\beta}=2\frac{\langle\beta,\alpha\rangle}{\langle\alpha,\alpha\rangle},
\qquad\alpha,\beta\in R.
\]
这称为 $\g$ 的约化根系 $R$ 的 Killing 形式。

元素 $\gamma\in\mathfrak{t}^{*}$ 称为正则的，若对所有 $\alpha\in R$ 有 $\langle\alpha,\gamma\rangle\ne0$。对正则元素 $\gamma$，令
\[
R^{+}(\gamma):=\{\alpha\in R\mid\langle\alpha,\gamma\rangle>0\}.
\]
根 $\alpha\in R^{+}(\gamma)$ 称为单的，若它不能写成 $R^{+}(\gamma)$ 中两个元素之和。单根确实构成 $R$ 的一组基。若根系不能分解为两个互不相交、非空且正交的子集之并，则称该根系是不可约的；这里的正交性按 Killing 形式理解。

% [原文 p.427]
\subsection{约化根系的分类}

我们现在看如何从一类称为 Dynkin 图的带权图确定复单李代数的根系，并且事实上在同构意义下完全分类所有这类李代数。这些图确定一个整数矩阵，称为 Cartan 矩阵，后者又确定根系。

\begin{definition}[A.15] 设 $\g$ 是秩为 $l$ 的复单李代数。设 $R$ 是 $\mathfrak{t}^{*}$ 的约化根系，$B=\{\alpha_1,\dots,\alpha_n\}$ 是 $R$ 的一组基。矩阵 $(A_{ij})_{1\le i,j\le l}$，其中 $A_{ij}=a_{\alpha_i,\alpha_j}$，称为 $\g$ 的 Cartan 矩阵；回忆 $a_{\alpha_i,\alpha_j}=2\frac{\langle\alpha_i,\alpha_j\rangle}{\langle\alpha_i,\alpha_i\rangle}$，对 $\alpha_i,\alpha_j\in R$。
\end{definition}

现在设 $R$ 是秩为 $l$ 的复单李代数 $\g$ 的 Cartan 子代数 $\mathfrak{t}$ 的不可约约化根系，$R$ 的基为 $B=\{\alpha_1,\dots,\alpha_l\}$，并设 $R^{+}$ 是对应的正根集。$R$ 的 Coxeter 图是有 $l$ 个顶点的图，第 $i$ 个顶点与第 $j$ 个顶点由 $a_{ij}a_{ji}$ 条互不相交的线相连，其中 $a_{ij}=a_{\alpha_i,\alpha_j}$。$R$ 的不可约性保证 Weyl 群 $W(R)$ 在 $\mathfrak{t}^{*}$ 上不可约地作用，从而不变内积 $\langle,\rangle$ 在相差常数因子的意义下唯一。在 Coxeter 图的第 $i$ 个顶点上赋予与 $\langle\alpha_i,\alpha_i\rangle$ 成比例的权，所得的图称为 Dynkin 图。在 Killing 与 Cartan 的开创性工作中，复单李代数的完全分类正是以 Dynkin 图的形式给出的。下面我们列出这些图，并给出不可约约化根系中一组基里的根，因为这在第 12 章中会用到；详细论述请读者参见 [195, 第 10 章, 第 3.3 节]。

% [Tutor Note: 定义 A.15 中基写作 $B=\{\alpha_1,\dots,\alpha_n\}$，而矩阵下标却用 $1\le i,j\le l$（$l$ 为秩）；同理 A.1.1 节中单根集写作 $\Pi=\{\alpha_1,\dots,\alpha_n\}$ 而 $n$ 是秩。原文 $n$ 与 $l$ 混用，实际都指秩；译文中 A.1.1 处照录 $n$，此处照录 $n$ 与 $l$ 的混用，请读者按「基有 $l$ 个元素」理解。]

\begin{theorem}[A.16] 秩为 $l$ 的复单李代数的 Cartan 子代数的约化根系的 Dynkin 图分类如下：
\begin{itemize}
\item 对李代数 $\mathfrak{a}_l\cong\sl_{l+1}(\C)$（$l\ge1$），Dynkin 图是一条单链 $\alpha_1,\alpha_2,\dots,\alpha_{l-1},\alpha_l$（相邻顶点以单边相连；顶点权全为 $1$），而约化根系的一组基由根
\[
\alpha_i=e_i-e_{i+1},\qquad 1\le i\le l
\]
给出。
\item 对李代数 $\mathfrak{b}_l\cong\mathfrak{so}_{2l+1}(\C)$（$l\ge2$），Dynkin 图是链 $\alpha_1,\alpha_2,\dots,\alpha_{l-1},\alpha_l$：前 $l-2$ 条边为单边，$\alpha_{l-1}$ 与 $\alpha_l$ 之间为双重边（箭头指向短根 $\alpha_l$）；顶点权为 $2,2,\dots,2,1$。约化根系的一组基由根
\[
\alpha_i=e_i-e_{i+1}\ (1\le i\le l-1),\qquad \alpha_l=e_l
\]
给出。
\item 对李代数 $\mathfrak{c}_l\cong\mathfrak{sp}_{2l}(\C)$（$l\ge3$），Dynkin 图是链 $\alpha_1,\alpha_2,\dots,\alpha_{l-1},\alpha_l$：前 $l-2$ 条边为单边，$\alpha_{l-1}$ 与 $\alpha_l$ 之间为双重边（箭头指向短根 $\alpha_{l-1}$）；顶点权为 $1,1,\dots,1,2$。约化根系的一组基由根
\[
\alpha_i=e_i-e_{i+1}\ (1\le i\le l-1),\qquad \alpha_l=2e_l
\]
给出。
% [原文 p.428]
\item 对李代数 $\mathfrak{d}_l\cong\mathfrak{so}_{2l}(\C)$（$l\ge4$），Dynkin 图在 $\alpha_{l-2}$ 处呈三叉形：水平链 $\alpha_1,\alpha_2,\dots,\alpha_{l-2}$（相邻以单边相连），而 $\alpha_{l-2}$ 又分别与 $\alpha_{l-1}$ 和 $\alpha_l$ 以单边相连；顶点权全为 $1$。约化根系的一组基由根
\[
\alpha_i=e_i-e_{i+1}\ (1\le i\le l-1),\qquad \alpha_l=e_{l-1}+e_l
\]
给出。
\item 对李代数 $\mathfrak{e}_6$，Dynkin 图是水平链 $\alpha_6,\alpha_5,\alpha_4,\alpha_3,\alpha_1$（相邻以单边相连），并在 $\alpha_4$ 处向上伸出一条单边连到 $\alpha_2$；顶点权全为 $1$。约化根系的一组基由根
\[
\alpha_1=\tfrac12(e_1+e_6)-\tfrac12(e_2+e_3+e_4+e_5),\quad \alpha_2=e_1+e_2,\quad \alpha_3=e_2-e_1,
\]
\[
\alpha_4=e_3-e_2,\quad \alpha_5=e_4-e_3,\quad \alpha_6=e_5-e_4
\]
给出。
\item 对李代数 $\mathfrak{e}_7$，Dynkin 图是水平链 $\alpha_7,\alpha_6,\alpha_5,\alpha_4,\alpha_3,\alpha_1$（相邻以单边相连），并在 $\alpha_5$ 处向上伸出一条单边连到 $\alpha_2$；顶点权全为 $1$。约化根系的一组基由根
\[
\alpha_1=\tfrac12(e_1+e_7)-\tfrac12(e_2+e_3+e_4+e_5+e_6),\quad \alpha_2=e_1+e_2,\quad \alpha_3=e_2-e_1,
\]
\[
\alpha_4=e_3-e_2,\quad \alpha_5=e_4-e_3,\quad \alpha_6=e_5-e_4,\quad \alpha_7=e_6-e_5
\]
给出。
\item 对李代数 $\mathfrak{e}_8$，Dynkin 图是水平链 $\alpha_8,\alpha_7,\alpha_6,\alpha_5,\alpha_4,\alpha_3,\alpha_1$（相邻以单边相连），并在 $\alpha_6$ 处向上伸出一条单边连到 $\alpha_2$；顶点权全为 $1$。约化根系的一组基由根
\[
\alpha_1=\tfrac12(e_1+e_8)-\tfrac12(e_2+e_3+e_4+e_5+e_6+e_7),\quad \alpha_2=e_1+e_2,\quad \alpha_3=e_2-e_1,
\]
\[
\alpha_4=e_3-e_2,\quad \alpha_5=e_4-e_3,\quad \alpha_6=e_5-e_4,\quad \alpha_7=e_6-e_5,\quad \alpha_8=e_7-e_6
\]
给出。
\item 对李代数 $\mathfrak{f}_4$，Dynkin 图是链 $\alpha_1,\alpha_2,\alpha_3,\alpha_4$：$\alpha_1$ 与 $\alpha_2$ 之间为单边，$\alpha_2$ 与 $\alpha_3$ 之间为双重边（箭头指向短根 $\alpha_3$），$\alpha_3$ 与 $\alpha_4$ 之间为单边；顶点权为 $2,2,1,1$。约化根系的一组基由根
\[
\alpha_1=e_2-e_3,\quad \alpha_2=e_3-e_4,\quad \alpha_3=e_4,\quad \alpha_4=\tfrac12(e_1-e_2-e_3-e_4)
\]
给出。
% [原文 p.429]
\item 对李代数 $\mathfrak{g}_2$，Dynkin 图由两个顶点 $\alpha_2$ 与 $\alpha_1$ 以三重边相连构成（箭头指向短根 $\alpha_1$）；顶点权为 $3,1$。约化根系的一组基由根
\[
\alpha_1=e_1-e_2,\qquad \alpha_2=-2e_1+e_2+e_3
\]
给出。
\end{itemize}
\end{theorem}

% [Tutor Note: 原书此处以 9 张 Dynkin 图排版给出；译文以链结构、边重数与顶点权的文字形式等价给出，信息完整。]

\section{实李群}

复半单李代数的结构理论具有若干良好性质，一般比实李代数简单。然而为我们的目的，我们需要考虑实李代数，而这些可以通过复李代数的实形式引入。半单李代数结构理论详细处理的标准参考文献包括 [195, 234, 291]。

\begin{definition}[A.17] 设 $\g_\R$ 是实李代数。复李代数 $\g_\C=\g_\R\otimes\C$ 称为 $\g_\R$ 的复化，而 $\g_\R$ 称为 $\g_\C$ 的实形式。
\end{definition}

\begin{remark}[A.18] 复化 $\g_\C$ 是单的（相应地半单的）当且仅当 $\g_\R$ 是单的（相应地半单的）。
\end{remark}

\begin{example}[A.19] $\mathfrak{sl}(n,\R)$ 与 $\mathfrak{su}(n)$ 的复化都是 $\mathfrak{sl}(n,\C)$，这表明两个不同构的实李代数可以有相同的复化。
\end{example}

\begin{proposition}[A.20] 设 $\g_\C$ 是复李代数，$\sigma:\g_\C\to\g_\C$ 是复反线性对合。对合 $\sigma$ 把 $\g_\C$ 分解为其 $(\pm1)$-特征空间：
\[
\g_\C=\g_+\oplus\g_-.
\]
子代数 $\g_+$ 是 $\g_\C$ 的一个实形式。反之，$\g_\C$ 的任一实形式 $\g_\R$ 都是同构 $\g_\C=\g_\R\otimes\C$ 在 $\g_\C$ 上诱导的复共轭的 $(+1)$-特征空间。
\end{proposition}

当 $\g_\C$ 半单时，其 Killing 形式 $K_\C$ 非退化，并限制为任一实形式 $\g_\R$ 上的非退化 Killing 形式 $K_\R$。由 Sylvester 惯性定律，实 Killing 形式 $K_\R$ 可在适当的基下对角化，对角元为 $+1$ 或 $-1$。

\begin{definition}[A.21] 设 $\g_\R$ 是实半单李代数。Cartan 对合是李代数对合 $\theta:\g_\R\to\g_\R$，使得双线性形式
\[
K_\theta:\g_\R\times\g_\R\to\R,\qquad (X,Y)\mapsto K_\R(X,-\theta(Y))
\]
正定。Cartan 对合 $\theta$ 把 $\g_\R$ 作为向量空间分解为其 $(-1)$-特征空间 $\mathfrak{p}$ 与 $(+1)$-特征空间 $\mathfrak{k}$，称为 Cartan 分解
\[
\g_\R=\mathfrak{p}\oplus\mathfrak{k}.
\]
\end{definition}

% [原文 p.430]
实半单李代数 $\g_\R$ 可能有多个 Cartan 对合，但它们都由 $\g_\R$ 的自同构共轭。注意 $(+1)$-特征空间 $\mathfrak{k}$ 是 $\g_\R$ 的子代数，而 $(-1)$-特征空间 $\mathfrak{p}$ 只是向量空间，因为 $[\mathfrak{p},\mathfrak{p}]\subset\mathfrak{k}$。因此，若实李群 $G$ 的李代数等于 $\g_\R$，则存在子群 $K\subset G$，其李代数为 $\mathfrak{k}$。该子群是 $G$ 的极大紧子群。还要注意 Killing 形式 $K_\R$ 在 $\mathfrak{p}$ 上正定、在 $\mathfrak{k}$ 上负定。

实半单李代数也允许权空间分解。我们称 $\mathfrak{p}$ 的极大阿贝尔子空间 $\mathfrak{a}$ 为 Cartan 子空间。于是李代数 $\g_\R$ 分解为
\[
\g_\R=\g_0\oplus\bigoplus_{\alpha\in R}\g_\alpha,
\]
其中 $R\subset\mathfrak{a}^{*}$ 是 $\g_\R$ 的限制根系。对每个 $\alpha$，空间 $\g_\alpha$ 是限制根 $\alpha\in\mathfrak{a}^{*}$ 的特征空间。与复情形不同，一般有 $\mathfrak{a}\ne\g_0$。$\mathfrak{a}$ 的维数称为 $\g_\R$ 的实秩，记作 $\operatorname{rk}_\R(\g_\R)$。

在复半单李代数的所有实形式中，有两个特定的实形式是我们所关心的：分裂实形式与紧形式。我们先描述紧形式。

\begin{definition}[A.22] $\g_\C$ 的实形式 $\g_\R$ 称为紧形式，若 $K_\R$ 在 $\g_\R$ 上负定。
\end{definition}

紧形式这个名称由下述命题说明其合理性：

\begin{proposition}[A.23] 实单李群是紧的当且仅当其李代数是紧形式。
\end{proposition}

\begin{remark}[A.24] 每个复半单李代数都有唯一的（模内自同构）紧形式。在此情形下，存在唯一的 Cartan 对合，即恒等映射。
\end{remark}

我们现在描述分裂实形式。

\begin{definition}[A.25] 对带 Cartan 子代数 $\mathfrak{t}$ 的复半单李代数 $\g_\C$，若其实形式 $\g_\R$ 的 Cartan 子空间 $\mathfrak{a}$ 满足 $\mathfrak{a}\otimes\C=\mathfrak{t}$，或等价地 $\operatorname{rk}_\C(\g_\C)=\operatorname{rk}_\R(\g_\R)$，则称 $\g_\R$ 为 $\g_\C$ 的分裂实形式。
\end{definition}

任何复半单李代数都含有分裂实形式；由单李代数的分类，我们可以列出分裂单实李代数在同构意义下的清单。

我们现在考虑这些李代数的指数以过渡到李群，从而定义：

\begin{definition}[A.26] 若 $\operatorname{rk}_\R G=\operatorname{rk}_\C G^{\C}$，则称单实李群 $G$ 为分裂实李群，其中 $G^{\C}:=G\otimes\C$ 表示 $G$ 的复化。
\end{definition}

% [原文 p.431]
\begin{fact}[A.27] 在相差覆叠的意义下，单分裂实李群是下列群之一：
\[
\SL_n(\R),\ \SO(n,n+1),\ \mathrm{Sp}_{2n}(\R),\ \SO(n,n),\ E_{6(6)},\ E_{7(7)},\ E_{8(8)},\ F_{4(4)},\ G_2,
\]
其中 $n\in\N$。前四个群称为经典群，后五个群称为例外群。
\end{fact}

对紧单李群有类似的分类结果：

\begin{fact}[A.28] 在相差覆叠的意义下，紧单李群是下列群之一：
\[
\SU(n),\ \SO(2n+1),\ \mathrm{USp}(2n),\ \SO(2n),\ E_6^c,\ E_7^c,\ E_8^c,\ F_4^c,\ G_2^c,
\]
其中 $n\in\N$。前四个群称为经典群，后五个群称为例外群。注意 $\mathrm{USp}(2n)=\mathrm{Sp}_{2n}(\C)\cap\SU(2n)$。
\end{fact}

% [Tutor Note: 事实 A.27 中 $E_{6(6)}$ 等是 Tits 记号：括号内数字为该实形式的实秩。按定义 A.25 与 A.26，「分裂」恰要求实秩等于复秩，故 $E_{6(6)}$、$E_{7(7)}$、$E_{8(8)}$、$F_{4(4)}$、$G_2$（即 $G_{2(2)}$）正是例外单李群的分裂实形式；而 $E_{6(2)}$、$E_{7(-5)}$ 等非分裂实形式不在该清单中。]

\section{Levi 分解、抛物子代数与旗簇}

在某些语境下，我们需要考虑非约化的李代数，例如 $\mathfrak{gl}_n$ 中由上三角矩阵构成的子代数。我们在此简要回顾 Levi 分解、Borel 子代数、抛物子代数与部分旗簇等概念；这些主题的标准参考文献是 Helgason 的书 [195]。

本节中，我们用 $\g$ 表示特征零域（通常是 $\R$ 或 $\C$）上的有限维李代数。回忆其根基 $\operatorname{rad}\g$ 是极大的可解理想。

\begin{theorem}[A.29]（Levi 分解）特征零域上任何有限维李代数都是其根基与一个半单李子代数的半直积：
\[
\g\cong\operatorname{rad}\g\rtimes\mathfrak{l}.
\]
半单子代数 $\mathfrak{l}$ 称为 Levi 子代数。
\end{theorem}

\begin{theorem}[A.30]（Levi–Malcev）任意两个 Levi 子代数都由内自同构共轭。
\end{theorem}

这一分解使人们可以把半单与可解李代数分开处理。我们通常把它应用于 $\mathfrak{gl}_n$ 中上三角矩阵的推广。

\begin{definition}[A.31] Borel 子代数 $\mathfrak{b}$ 是 $\g$ 的极大可解子代数。其换位子 $\mathfrak{u}=[\mathfrak{b},\mathfrak{b}]$ 称为幂幺子代数。
\end{definition}

\begin{example}[A.32] 对 $\g=\mathfrak{gl}_n$，上三角矩阵的空间构成 Borel 子代数，而严格上三角矩阵的空间构成幂幺子代数。
\end{example}

% [原文 p.432]
\begin{example}[A.33] 更一般地，对 $\g=\mathfrak{gl}(V)$，其中 $V$ 是 $n$ 维有限维向量空间，Borel 子代数是完全旗的稳定化子。给定
\[
V_0=\{0\}\subset V_1\subset V_2\subset\dots\subset V_n=V,\qquad \dim(V_i)=i,
\]
则集合 $\{M\in\mathfrak{gl}(V)\mid M(V_i)\subset V_i\ \forall i\}$ 是 Borel 子代数。$\mathfrak{gl}(V)$ 的任何 Borel 子代数都是这种形式。
\end{example}

\begin{proposition}[A.34] 任意两个 Borel 子代数都由内自同构共轭。
\end{proposition}

李代数的一类更一般的子代数是那些包含 Borel 子代数的子代数：

\begin{definition}[A.35] 包含 Borel 子代数的真子代数 $\mathfrak{p}\subset\g$ 称为抛物的（parabolic）。
\end{definition}

\begin{example}[A.36] 对如前所述的 $\g=\mathfrak{gl}(V)$，抛物子代数是部分旗的稳定化子。给定部分旗 $V_{i_0}\subset V_{i_1}\subset\dots\subset V_{i_m}$（其中 $0\le i_1<i_2<\dots<i_m\le n$），空间 $\{M\in\mathfrak{gl}(V)\mid M(V_{i_k})\subset V_{i_k}\}$ 是抛物子代数。它包含与部分旗的任一补全相联系的那个 Borel 子代数。$\mathfrak{gl}(V)$ 中的任何抛物子代数都是这种形式。
\end{example}

一般地，抛物子代数与单根集 $\Delta$ 的非空子集相对应：

\begin{proposition}[A.37] 模去内自同构，抛物子代数与非空子集 $\Theta\subset\Delta$ 之间存在双射。
\end{proposition}

特别地，共有 $2^{\operatorname{rk}(\g)}-1$ 种不同的抛物子代数。若 $|\Theta|=1$，则抛物子代数是极大的；当 $\Theta=\Delta$ 时，抛物子代数极小，且等于 Borel 子代数。

对给定的子集 $\Theta\subset\Delta$，抛物子代数 $\mathfrak{p}_\Theta$ 是由所有 $f_i$、$h_i$ 以及那些满足「对所有 $\beta\in\Theta$ 有 $\beta(e_i)\ne0$」的 $e_i$ 生成的李代数，其中 $(f_i,h_i,e_i)$ 表示 $\g$ 的一组 Chevalley 基。对立的抛物子代数 $\mathfrak{p}_\Theta^{\mathrm{opp}}$ 由所有 $e_i$、$h_i$ 以及那些满足「对所有 $\beta\in\Theta$ 有 $\beta(f_i)\ne0$」的 $f_i$ 生成。

把 Levi 分解应用于抛物子代数，给出：

\begin{proposition}[A.38] 抛物子代数是其幂幺根基与一个半单子代数的半直积。
\end{proposition}

\begin{example}[A.39] 考虑带内积的向量空间 $V$。对 $\g=\mathfrak{gl}(V)$ 与对应于部分旗 $(V_{i_k})_{1\le k\le m}$ 的抛物子代数 $\mathfrak{p}$，$V^{*}$ 中对立的旗可以等同于正交补 $(V_{i_k}^{\perp})$。此时 Levi 子代数是交集 $\mathfrak{p}\cap\mathfrak{p}^{\mathrm{opp}}$，而幂幺根基是同时保持 $\mathfrak{p}$ 与 $\mathfrak{p}^{\mathrm{opp}}$ 的变换构成的空间。
\end{example}

在李群层面，我们有类似的概念。设 $G$ 是李群，其李代数为 $\g$。

% [原文 p.433]
\Fig[0.85\linewidth]{figures/fig_A_1.png}{A.1}{Levi 子群、抛物子群及其幂幺根（以 $\GL_n$ 为例）的示意。}

\begin{definition}[A.40] Borel 子群 $B$ 是 $G$ 的极大可解子群。其换位子群 $U=[B,B]$ 是幂幺子群。抛物子群 $P$ 是 $G$ 中包含 Borel 子群的真正子群。
\end{definition}

李群与其李代数之间的一般对应给出 Borel 子群与 Borel 子代数之间的双射，抛物子群与幂幺子群同相应子代数之间也是如此。如前所述，抛物子群与部分旗之间有密切关系：前者是后者的稳定化子。

\begin{definition}[A.41] 李群 $G$ 的旗簇是陪集空间 $G/B$，其中 $B$ 是 Borel 子群。部分旗簇是陪集空间 $G/P$，其中 $P$ 是 $G$ 的抛物子代数。
\end{definition}

% [Tutor Note: 定义 A.41 原文把 $P$ 称为「$G$ 的抛物子代数」（parabolic subalgebra），但 $G/P$ 是李群的商、$P$ 应为抛物子群；这与定义 A.40 中「抛物子群」的定义不一致，属原文笔误。译文照录并在此指出。]

\begin{example}[A.42] 对 $G=\SL_2(\C)$，旗的空间可与 $\mathbb{CP}^1$ 等同。$G$ 在完全旗上的作用可迁，且稳定化子是 Borel 子群，因此有 $G/B\cong\mathbb{CP}^1$。
\end{example}

更一般地，对 $G=\mathrm{GL}(V)$，$G$ 在所有完全旗的空间上的作用可迁。固定旗的稳定化子是 Borel 子群。因此，我们可以把 $G/B$ 与所有完全旗的空间等同。类似地，$G/P$ 可与所有部分旗的空间等同，其类型由 $P$ 的稳定化子指定。

在存在内积时，正交化过程（或 Gram–Schmidt 算法）使我们能把一组标准正交基联系到一个完全旗。下面的命题可以看作这一点的推广。

\begin{proposition}[A.43] 旗簇 $G/B$ 微分同胚于陪集空间 $K/T$，其中 $K$ 是 $G$ 的极大紧子群，$T\subset K$ 是 $K$ 中的极大环面。更一般地，对部分旗簇有 $G/P\cong K/L_\R$，其中 $L_\R=L\cap K$，而 $L$ 是 $P$ 对应的 Levi 子群。
\end{proposition}

\begin{example}[A.44] 把命题应用于例 A.42，得
\[
\mathbb{CP}^1\cong\SL_2(\C)/B\cong\SU(2)/T\cong\SO(3)/\SO(2).
\]
\end{example}

% [原文 p.434]
% 原书此页为空白页。

% [原文 p.435]
% 原书此页为 Glossary 部分起始的空白页（仅保留页眉「Glossary 435」，无正文）。

% [原文 p.436]
% 原书此页为空白页。

% [Tutor Note: 例 A.42 与例 A.44 的链条可作维数核对：$\SL_2(\C)$ 复维 3，其 Borel 子群 $B$ 复维 2，故 $\SL_2(\C)/B$ 复维 1、即实维 2；$\SU(2)$ 实维 3、$T$ 实维 1，商 $\SU(2)/T$ 实维 2；$\SO(3)$ 实维 3、$\SO(2)$ 实维 1，商 $\SO(3)/\SO(2)$ 实维 2；三者都与 $\mathbb{CP}^1$ 的实维 2 一致。]
